% arXiv preprint. Build: pdflatex main && pdflatex main
\documentclass[10pt,a4paper]{article}
\usepackage[T1]{fontenc}
\usepackage[utf8]{inputenc}
\usepackage{lmodern}
\usepackage[french,english]{babel}
\usepackage[a4paper,margin=0.92in,top=1in,bottom=1in]{geometry}
\usepackage{microtype}
\usepackage{amsmath,amssymb}
\usepackage{booktabs}
\usepackage{graphicx}
\usepackage{xcolor}
\usepackage{fancyhdr}
\usepackage{titlesec}
\usepackage{caption}
\usepackage{enumitem}
\usepackage{tikz}
\usepackage{pgfplots}
\pgfplotsset{compat=1.17}
\usetikzlibrary{arrows.meta,positioning,shapes.geometric,fit,calc,backgrounds}
\usepackage[hidelinks]{hyperref}

% ---- author block: edit here -------------------------------------------------
\newcommand{\paperauthor}{Joey Lukoji}
\newcommand{\paperaffil}{AURIX Project}
\newcommand{\papercontact}{\url{https://github.com/joeylukoji/Aurix}}
% ------------------------------------------------------------------------------

\definecolor{adv}{HTML}{2F6DB5}
\definecolor{dec}{HTML}{1E8A5A}
\definecolor{cus}{HTML}{B5642F}
\definecolor{ink}{HTML}{222222}
\definecolor{mute}{HTML}{6B6B6B}

\pagestyle{fancy}\fancyhf{}
\fancyhead[L]{\small\textsc{Verifiable Conditional Payments}}
\fancyhead[R]{\small\textsc{A Preprint}}
\fancyfoot[C]{\small\thepage}
\renewcommand{\headrulewidth}{0.4pt}
\titleformat{\section}{\large\bfseries}{\thesection}{0.6em}{}
\titleformat{\subsection}{\normalsize\bfseries}{\thesubsection}{0.6em}{}
\titlespacing*{\section}{0pt}{9pt}{4pt}
\titlespacing*{\subsection}{0pt}{7pt}{3pt}
\captionsetup{font=small,labelfont=bf,skip=5pt}
\setlength{\parskip}{2.5pt plus 1pt}
\setlength{\textfloatsep}{9pt plus 2pt}
\setlength{\floatsep}{8pt plus 2pt}
\setlist{itemsep=1pt,topsep=2pt,parsep=0pt,leftmargin=1.4em}
\newcommand{\code}[1]{\texttt{\small #1}}
\newcommand{\rel}{\textsc{release}}
\newcommand{\void}{\textsc{void}}
\newcommand{\hold}{\textsc{hold}}
\newcommand{\wait}{\textsc{wait}}

\begin{document}
\thispagestyle{plain}

\begin{center}
\rule{\textwidth}{1.6pt}\\[9pt]
{\Large\bfseries A General Methodology for AI-Assisted, Evidence-Based and\\[3pt]
Cryptographically Verifiable Conditional Payment Systems\par}
\vspace{6pt}\rule{\textwidth}{0.5pt}\\[10pt]
{\large \paperauthor}\\[2pt]
{\small \paperaffil\\ \papercontact}\\[4pt]
{\small November 2026}
\end{center}

\begin{abstract}
\noindent A conditional payment moves money only if an agreed event took place. Card and wallet providers
implement the condition as a manual capture by the merchant, and blockchains implement it as contract code fed by an
oracle. In both cases a third party who was not present cannot check afterwards why the money moved. We describe a
methodology that separates such a system into three planes with different authority. A custody plane holds the
funds and is the only one that can move them. A decision plane is a pure, versioned function of signed evidence and
of the time at which each item was received. An advisory plane, where the machine learning lives, drafts conditions
and scores risk but has no write access to the decision. Each decision is exported with its inputs as a signed
package, and an independent program replays the function offline and must reach the same outcome. We give the
evidence model, the semantics of a small condition language, the decision function and the rules that keep model
output outside the trusted path. A reference implementation, AURIX, runs the method over a commercial payment
provider and a four-validator BFT ledger. On the PaySim benchmark, with a time-ordered split and thresholds fixed on
validation data, a calibrated logistic regression reaches a test PR-AUC of 0.973 (95\% CI 0.967--0.978), while a
gradient-boosted model is rejected because its validation PR-AUC ranges from 0.004 to 0.81 across five seeds.
Exhaustive single-bit corruption of three proof packages (192,880 trials) never altered terms, evidence or a decision
without detection, but it exposed two defects that we report: 123 flips in the unsigned envelope of the manifest
signature were accepted, and 1,348 trials ended in an unhandled exception instead of a report. We state what the
method does not establish: a valid signature says who asserted a delivery, not that the goods arrived.
\end{abstract}

\begin{otherlanguage}{french}
\begin{abstract}
\noindent Un paiement conditionnel ne déplace les fonds que si un événement convenu a eu lieu. Les prestataires de
paiement traitent la condition comme une capture manuelle par le marchand, et les blockchains comme du code de
contrat alimenté par un oracle. Dans les deux cas, un tiers absent au moment des faits ne peut pas vérifier après
coup pourquoi l'argent a bougé. Nous décrivons une méthodologie qui sépare un tel système en trois plans d'autorité
distincte. Le plan de garde détient les fonds et lui seul peut les déplacer. Le plan de décision est une fonction
pure et versionnée des preuves signées et de l'heure de réception de chacune. Le plan consultatif, où réside
l'apprentissage automatique, rédige les conditions et évalue le risque, sans aucun accès en écriture à la décision.
Chaque décision est exportée avec ses entrées dans un paquet signé, et un programme indépendant rejoue la fonction
hors ligne et doit aboutir au même résultat. Nous donnons le modèle de preuve, la sémantique d'un petit langage de
conditions, la fonction de décision et les règles qui maintiennent la sortie des modèles hors du chemin de confiance.
Une implémentation de référence, AURIX, applique la méthode au-dessus d'un prestataire de paiement commercial et d'un
registre BFT à quatre validateurs. Sur le jeu PaySim, avec un découpage chronologique et des seuils fixés sur les
données de validation, une régression logistique calibrée atteint une PR-AUC de test de 0,973 (IC à 95\,\% :
0,967--0,978), tandis qu'un modèle de gradient boosting est écarté car sa PR-AUC de validation varie de 0,004 à 0,81
selon la graine. La corruption exhaustive, bit par bit, de trois paquets de preuve (192\,880 essais) n'a jamais
modifié les termes, les preuves ou une décision sans être détectée, mais elle a révélé deux défauts que nous
signalons : 123 inversions dans l'enveloppe non signée de la signature du manifeste ont été acceptées, et
1\,348 essais se sont terminés par une exception non gérée au lieu d'un rapport. Nous précisons ce que la méthode n'établit pas : une signature valide indique qui a
affirmé une livraison, non que la marchandise est arrivée.
\end{abstract}
\end{otherlanguage}

\noindent{\small\textbf{Keywords:} conditional payment, escrow, signed attestation, deterministic decision rule,
offline verification, fraud scoring, calibration, large language models, BFT ledger.\\
\textbf{Subjects:} cs.CR (primary); cs.AI; stat.ML; q-fin.GN.}

\section{Introduction}

A buyer in Kinshasa orders machine parts from a supplier in Guangzhou. The buyer does not want to pay before the
parts clear customs; the supplier does not want to ship before the money is committed. The usual answers are a
letter of credit, which costs a bank fee and several days, or an authorisation on a card or wallet that the
merchant captures by hand. The first is out of reach for small amounts. The second leaves the release in the
merchant's hands and gives the buyer a dispute procedure only after the charge.

Fair exchange between two parties who do not trust each other cannot be achieved without a third
party~\cite{pagnia1999}. Optimistic protocols limit the third party to the cases where something goes
wrong~\cite{asokan2000}. Smart contracts~\cite{szabo1997} replace it with a replicated program, and contract
languages such as the combinators of Peyton Jones et al.~\cite{peytonjones2000}, BitML~\cite{bartoletti2018} and
Marlowe~\cite{lamela2018} give that program a precise meaning. A contract on a ledger, however, sees only what is
written to the ledger. A delivery, an inspection or a customs clearance has to be reported by an
oracle~\cite{zhang2016,zhang2020}, and the ledger cannot tell a truthful report from a false
one~\cite{caldarelli2020}. Most consumer money also sits with regulated providers, not in on-chain escrow.

Machine learning adds a third difficulty. A language model can turn ``pay when the goods are delivered and
inspected, within ten days'' into a formal condition, and a classifier can flag a suspicious order. Neither output
is reproducible in the sense an auditor needs. The same prompt can give two answers, and text retrieved by the
model can carry instructions planted by an attacker~\cite{greshake2023}. If a model takes part in the release
decision, the decision can no longer be re-derived by someone else.

This paper gives a method for building conditional payment systems under those constraints. Its contributions are:
\begin{enumerate}
\item an architecture in three planes (custody, decision, advisory) with an explicit rule on which plane may write
what (Section~\ref{sec:method});
\item an evidence model and a decision function defined over signed attestations and their reception times, with a
condition language small enough to be replayed by an independent verifier;
\item an integration discipline for learned components: schema-bounded drafting approved by a human, and an
advisory risk score with pre-registered model selection, calibration and exact attributions;
\item a reference implementation and two evaluations whose raw outputs are in the repository: a fraud-scoring
study on PaySim and a bit-level tamper study of the proof format (Section~\ref{sec:eval}).
\end{enumerate}

\section{Related work}\label{sec:related}

\paragraph{Exchange protocols and contract languages.} Asokan, Shoup and Waidner~\cite{asokan2000} exchange
signatures with an arbiter that intervenes only on failure; our arbiter has the same optimistic role, restricted to
a hold and a final ruling. Contract languages~\cite{peytonjones2000,bartoletti2018,lamela2018} show that a small
set of combinators covers a wide range of financial agreements. We borrow that restraint: our conditions are trees
with three connectives, bounded in size, so that two independent implementations can be kept in agreement by tests.

\paragraph{Oracles.} Town Crier~\cite{zhang2016} authenticates web data with trusted hardware and
DECO~\cite{zhang2020} proves statements about a TLS session without server changes. Both strengthen the link
between a data source and a contract. They do not remove the source's ability to state something false, which
Caldarelli calls the oracle problem~\cite{caldarelli2020}. We take the same position and make the residual trust
explicit in every verification report.

\paragraph{Tamper evidence.} Linked time-stamping~\cite{haber1991} and transparency logs~\cite{rfc6962} let a
client detect that a record was altered or back-dated. We use a BFT ledger of the Tendermint
family~\cite{buchman2018} for the same purpose: it stores a commitment to the outcome and a quorum of confirmations,
not the business logic. Signatures are Ed25519~\cite{bernstein2012,rfc8032} over a canonical JSON
form~\cite{rfc8785}.

\paragraph{Fraud scoring.} Operational fraud detection faces class imbalance, drift and delayed
labels~\cite{dalpozzolo2018}. Under imbalance, precision--recall analysis is more informative than
ROC~\cite{davis2006,saito2015}. Raw classifier scores are poor probabilities and need
calibration~\cite{zadrozny2002,niculescu2005}. Kapoor and Narayanan~\cite{kapoor2023} trace many inflated results to
leakage between training and evaluation data, and preregistration~\cite{nosek2018} is the standard defence against
choices made after seeing the outcome. Section~\ref{sec:risk} applies these points and reports one deviation from
our own plan.

\section{Methodology}\label{sec:method}

\subsection{Three planes and one write rule}

Figure~\ref{fig:arch} shows the decomposition. The \emph{custody plane} is whoever legally holds the money: a
payment provider that supports authorise, capture and void, or an on-chain escrow account. The \emph{decision
plane} computes one of four outcomes from signed inputs. The \emph{advisory plane} contains every component whose
output is statistical: the language model and the risk scorer.

The write rule is: custody acts only on a signed decision; the decision reads only signed terms and signed
evidence; advisory output reaches the decision only by becoming terms that a human has read and approved. The
project states it as \emph{the AI advises, the rule decides, the human approves}. The operator of the decision
plane never holds the funds, which keeps the method usable on top of an existing regulated provider.

\begin{figure}[t]
\centering
\resizebox{\textwidth}{!}{%
\begin{tikzpicture}[font=\footnotesize, node distance=4mm and 9mm,
  box/.style={draw=#1, fill=#1!8, rounded corners=2pt, minimum height=7.5mm, minimum width=24mm, align=center, text=ink},
  ext/.style={draw=mute, rounded corners=2pt, minimum height=7.5mm, minimum width=22mm, align=center, text=ink},
  arr/.style={-{Stealth[length=4.5pt]}, semithick, draw=ink},
  dsh/.style={-{Stealth[length=4.5pt]}, semithick, dashed, draw=adv},
  lab/.style={font=\scriptsize, text=mute, inner sep=1pt}]
% advisory
\node[box=adv] (llm) {Language model\\drafts a condition};
\node[box=adv, right=of llm] (val) {Schema validator\\bounds, known sources};
\node[box=adv, right=18mm of val] (risk) {Risk scorer\\calibrated, attributed};
% human
\node[ext, below=9mm of val] (buyer) {Buyer\\reads, edits, approves};
% decision
\node[box=dec, below=9mm of buyer] (terms) {Signed terms $T$};
\node[box=dec, right=of terms] (rule) {Decision function $\delta$\\pure, versioned};
\node[box=dec, right=of rule] (pkg) {Proof package $P$};
\node[box=dec, left=of terms] (ev) {Evidence store\\signature, role, receipt};
% sources
\node[ext, left=12mm of ev] (src) {Sources\\carrier, inspector,\\customs, arbiter};
% custody
\node[box=cus, below=9mm of rule] (psp) {Payment provider\\authorise / capture / void};
\node[box=cus, right=of psp] (chain) {BFT ledger\\$k$-of-$n$ attestors};
\node[ext, right=of pkg] (ver) {Independent\\verifier};
\draw[arr] (llm) -- (val);
\draw[dsh] (val) -- node[lab, right]{draft} (buyer);
\draw[dsh] (risk) |- node[lab, pos=0.25, right]{score, reasons} (buyer);
\draw[arr] (buyer) -- node[lab, right]{approval} (terms);
\draw[arr] (src) -- node[lab, above]{signed} (ev);
\draw[arr] (ev) -- (terms);
\draw[arr] (terms) -- (rule);
\draw[arr] (rule) -- (pkg);
\draw[arr] (rule) -- node[lab, right]{signed decision} (psp);
\draw[arr] (psp) -- (chain);
\draw[arr] (pkg) -- (ver);
\draw[arr] (chain) |- ($(pkg.south)+(0,-2.2mm)$) -- (pkg.south);
\begin{scope}[on background layer]
\node[fit=(llm)(risk), draw=adv!60, dashed, rounded corners=3pt, inner sep=5pt, label={[font=\scriptsize\bfseries, text=adv]above left:Advisory plane (no authority)}] {};
\node[fit=(ev)(pkg), draw=dec!60, dashed, rounded corners=3pt, inner sep=5pt, label={[font=\scriptsize\bfseries, text=dec]below left:Decision plane (deterministic)}] {};
\node[fit=(psp)(chain), draw=cus!60, dashed, rounded corners=3pt, inner sep=5pt, label={[font=\scriptsize\bfseries, text=cus]below left:Custody and anchoring}] {};
\end{scope}
\end{tikzpicture}}
\caption{The three planes. Solid arrows carry signed data. Dashed arrows carry model output, which stops at the
buyer: only the approved, signed terms enter the decision plane.}
\label{fig:arch}
\end{figure}

\subsection{Evidence}

An \emph{attestation} is a message $m$ stating that an event of a given type happened for a given order, together
with a signature under a registered key $k$:
\begin{equation}\label{eq:att}
a = (m, k, \sigma), \qquad \sigma = \mathrm{Sign}_{sk_k}\!\bigl(H(C(m))\bigr),
\end{equation}
where $C$ is a canonical serialisation~\cite{rfc8785} and $H$ a collision-resistant hash. A key registry gives each
key one role (the operator, a named source, the arbiter) and a validity interval. A source key may sign only the
event types of its role, so a carrier cannot attest an inspection.

The time written inside $m$ is chosen by the signer and can be back-dated. The operator therefore issues a signed
\emph{receipt} recording $r(a)$, the time at which it received $a$. For an order with terms $T$, deadline $D$ and
evaluation time $t$, the admissible evidence is
\begin{equation}\label{eq:adm}
E_t = \bigl\{\, a \;:\; \mathrm{Verify}_{pk_k}(a) \,\wedge\, \mathrm{role}(k) \ni \mathrm{type}(m) \,\wedge\,
\mathrm{ref}(m) = \mathrm{id}(T) \,\wedge\, k \text{ valid at } r(a) \,\wedge\, r(a) \le \min(t, D) \,\bigr\}.
\end{equation}
Evidence that arrived before the deadline but is evaluated after it still counts; evidence that arrived late does
not, whatever date it carries. Amounts are compared as decimals and references by exact equality.

\subsection{Conditions}

A condition $c$ is a tree. A leaf $\ell = (\tau, s, b)$ asks for an event of type $\tau$ from source $s$, optionally
dated before $b$. Internal nodes are \code{all}, \code{any} and \code{at\_least}~$k$. Satisfaction over an evidence
set $E$ is
\begin{equation}\label{eq:sem}
E \models c \iff
\begin{cases}
\exists\, a \in E:\ \mathrm{type}(a)=\tau,\ \mathrm{src}(a)=s,\ \mathrm{time}(a)\le b & c = (\tau,s,b),\\[2pt]
\bigl|\{\, i : E \models c_i \,\}\bigr| \ge k & c = \texttt{at\_least}_k(c_1,\dots,c_n),
\end{cases}
\end{equation}
with \code{all} and \code{any} the cases $k=n$ and $k=1$. Satisfaction is monotone: adding evidence never turns a
satisfied condition into an unsatisfied one. Relative limits such as ``within seven days'' are converted to absolute
dates when the terms are signed, so the signed terms never depend on when they are evaluated. A condition holds at
most ten leaves and four levels; the bound keeps a human able to read what they approve.
Figure~\ref{fig:fsm} (right) shows an example.

\subsection{The decision function}

Let $\mathit{auth}(E)$ hold when the custody plane has attested an authorisation for exactly the agreed amount and
currency, and let $t_{\mathrm{sat}}$ be the earliest reception time at which $E \models c$. A buyer may be granted a
dispute window of $w$ hours after the condition is met, so the earliest release time is
\begin{equation}\label{eq:trel}
t_{\mathrm{rel}} = \min\bigl(t_{\mathrm{sat}} + w,\; D\bigr).
\end{equation}
Write $\alpha$ for the latest signed arbiter action received before $D$, and $\mathit{disp}$ for an open buyer
dispute received before $t_{\mathrm{rel}}$. The decision is
\begin{equation}\label{eq:delta}
\delta(T, E_t, t) =
\begin{cases}
\alpha & \alpha \in \{\rel, \void\},\\
\void & t > D \text{ and not } \bigl(\mathit{auth}(E_t) \wedge E_t \models c \wedge \neg\mathit{disp} \wedge \alpha \ne \hold \bigr),\\
\hold & t \le D \text{ and } (\alpha = \hold \text{ or } \mathit{disp}),\\
\rel & \mathit{auth}(E_t) \wedge E_t \models c \wedge t \ge t_{\mathrm{rel}},\\
\wait & \text{otherwise.}
\end{cases}
\end{equation}
The cases are read top to bottom. $\delta$ takes no input other than $T$, $E_t$ and $t$: no clock read, no network
call, no floating-point arithmetic, no model. \rel{} and \void{} are final; the service signs the first final
outcome, and the capture or void sent to the custody plane carries an idempotency key derived from the order and
the decision, so a retry cannot produce the opposite action. Figure~\ref{fig:fsm} (left) gives the state diagram.

\begin{figure}[t]
\centering
\begin{minipage}[c]{0.53\textwidth}
\centering
\resizebox{\textwidth}{!}{%
\begin{tikzpicture}[font=\footnotesize, >={Stealth[length=4.5pt]},
  st/.style={draw=ink, rounded corners=8pt, minimum width=16mm, minimum height=7mm, fill=white},
  fin/.style={st, double, double distance=1pt},
  lab/.style={font=\scriptsize, text=mute, align=center, inner sep=2pt}]
\node[st] (wait) {\wait};
\node[st, right=24mm of wait] (hold) {\hold};
\node[fin, fill=dec!12, left=24mm of wait] (rel) {\rel};
\path (wait) -- coordinate (mid) (hold);
\node[fin, fill=cus!12, below=15mm of mid] (void) {\void};
\draw[<-] (wait.south west) -- ++(-6mm,-6mm) node[lab, below left, inner sep=1pt]{authorised};
\draw[->] (wait) to[bend left=14] node[lab, above]{dispute or\\arbiter hold} (hold);
\draw[->] (hold) to[bend left=14] node[lab, below]{unhold} (wait);
\draw[->] (wait) -- node[lab, above]{$E_t \models c$\\and $t \ge t_{\mathrm{rel}}$} (rel);
\draw[->] (wait) -- node[lab, anchor=east, pos=0.55, xshift=-4pt]{$t > D$} (void);
\draw[->] (hold) -- node[lab, anchor=west, pos=0.55, xshift=4pt]{$t > D$, or\\arbiter: void} (void);
\draw[->] (hold.north) to[out=140, in=40, looseness=0.75] node[lab, above]{arbiter: release} (rel.north);
\end{tikzpicture}}
\end{minipage}\hfill
\begin{minipage}[c]{0.43\textwidth}
\centering
\resizebox{\textwidth}{!}{%
\begin{tikzpicture}[font=\scriptsize,
  level 1/.style={sibling distance=30mm, level distance=12mm},
  level 2/.style={sibling distance=27mm, level distance=12mm},
  op/.style={draw=dec, fill=dec!10, rounded corners=2pt, inner sep=3pt},
  lf/.style={draw=mute, inner sep=3pt, align=center},
  edge from parent/.style={draw=ink, semithick}]
\node[op] {\texttt{all}}
  child {node[lf] {customs.cleared\\\textit{customs}}}
  child {node[lf] {shipment.delivered\\\textit{carrier}}}
  child {node[op] {\texttt{at\_least 1}}
    child {node[lf] {inspection.passed\\\textit{lab A}}}
    child {node[lf] {inspection.passed\\\textit{lab B}}}};
\end{tikzpicture}}
\end{minipage}
\caption{Left: states of an order under $\delta$. Double borders mark final states. A hold that is still in place
at the deadline ends in \void, so funds are never kept past the provider's authorisation period. Right: a condition
tree with four leaves and three levels.}
\label{fig:fsm}
\end{figure}

Three properties follow from the definitions, given unforgeable signatures and an honest key registry.
\emph{Safety:} no \rel{} is produced without a matching authorisation and either evidence satisfying the approved
condition or a signed arbiter ruling. \emph{Termination:} every order reaches \rel{} or \void{} at the first
evaluation after $D$ at the latest, since no branch returns \wait{} or \hold{} once $t > D$. \emph{Non-interference:}
advisory output does not appear among the arguments of $\delta$, so changing or removing the models changes no
decision for fixed terms and evidence.

\subsection{Proof packages and replay}

When an order reaches a final state the service exports a package $P$: the terms, each attestation with its
signature, the receipts, the decision $d$ with the hashes of the inputs it used and its evaluation time $t_d$, the
key registry, and a manifest listing the hash of every file. A verifier accepts the package only if all three
checks pass:
\begin{equation}\label{eq:verify}
\mathrm{Valid}(P) \iff \underbrace{\forall f \in P:\ H(f) = h_f}_{\text{integrity}} \;\wedge\;
\underbrace{\forall (x,\sigma) \in P:\ \mathrm{Verify}(x,\sigma)}_{\text{authenticity}} \;\wedge\;
\underbrace{\delta(T, E_{t_d}, t_d) = d}_{\text{replay}}.
\end{equation}
The third conjunct is what distinguishes the method from a signed receipt. The verifier does not trust the
operator's statement of the outcome; it recomputes the outcome from the inputs. An operator who captured without
sufficient evidence cannot produce a package that passes, short of forging a source signature. The verifier pins the
trusted keys out of band, since a package that ships its own registry would otherwise vouch for itself.

After a capture, $n$ independent attestors query the provider themselves, and the outcome is recorded on the ledger
once $k$ of them agree. This removes the operator as the single witness of the provider's state. It does not yet
replicate $\delta$ itself (Section~\ref{sec:limits}).

\subsection{Learned components}\label{sec:ai}

\paragraph{Drafting.} The language model receives the buyer's description and the list of registered sources, and
returns a candidate tree. The output is parsed against the grammar of Eq.~\eqref{eq:sem}, the size bounds and the
registry; anything else is discarded and a deterministic template is used, labelled as such. The summary and the
workflow diagram shown to the buyer are computed from the tree, never copied from model text, so the buyer approves
what will be executed and not what the model says it wrote. The model sees typed fields only (event types, sources,
amounts, dates), without payer identity or raw third-party payloads, which narrows the channel for injected
instructions~\cite{greshake2023}. Every generated text is stored in the package with the name of the model that
produced it.

\paragraph{Risk scoring.} The scorer returns a probability, a level and reasons before the buyer approves. It never
blocks. Candidates are compared on validation PR-AUC over $S$ training seeds, and a candidate is eligible only if
its spread across seeds is small:
\begin{equation}\label{eq:select}
m^\star = \arg\max_{m \in \mathcal{M}}\ \overline{\mathrm{PR}}_{\mathrm{val}}(m) \quad \text{subject to} \quad
\mathrm{sd}_{s=1..S}\,\mathrm{PR}_{\mathrm{val}}(m, s) \le \tau .
\end{equation}
The two alert thresholds have a fixed meaning, set on validation data only: the medium level must catch a stated
share of fraud, and the high level must be right a stated share of the time,
\begin{equation}\label{eq:thr}
\theta_{\mathrm{med}} = \max\{\theta : \mathrm{Recall}_{\mathrm{val}}(\theta) \ge \rho\}, \qquad
\theta_{\mathrm{high}} = \min\{\theta : \mathrm{Precision}_{\mathrm{val}}(\theta) \ge \pi\}.
\end{equation}
The displayed probability is the raw score mapped through isotonic regression fitted on the validation
split~\cite{zadrozny2002}. For a linear model on standardised features, the SHAP value~\cite{lundberg2017} of
feature $j$ has the closed form $\beta_j (x_j - \mu_j)/s_j$, so the reasons shown are exact and cost nothing to
compute. A quorum of two scorer instances signs each response.

\section{Reference implementation}\label{sec:impl}

AURIX\footnote{Source, tests and raw results: \url{https://github.com/joeylukoji/Aurix}. The token and the network
are a testnet without monetary value.} implements the method as a set of services behind a gateway. Custody is a
PayPal order with intent \code{AUTHORIZE}; the provider honours an authorisation for at most 29 days, which bounds
$D$. Evidence sources (carrier, inspector, customs) are simulated services, each with its own Ed25519 key, and are
flagged as simulated in the interface, the package and the verifier report. Two rule versions are deployed:
\code{delivery-attested/1} (one carrier attestation) and \code{conditions/1} (the tree language above). The
decision function is about 300 lines of Python without dependencies beyond the standard library and a signature
primitive. It is shared by the payment service and the command-line verifier and ported to JavaScript for an
in-browser verifier; a parity test runs both on the same packages. The ledger is a Cosmos SDK chain on CometBFT
with four validators and $k{=}2$ attestors. Drafting uses an open-weights model (Gemma) through an
OpenAI-compatible endpoint. Table~\ref{tab:compare} places the design next to the two common alternatives.

\begin{table}[t]
\centering\small
\caption{Where the release decision lives in three designs.}
\label{tab:compare}
\begin{tabular}{@{}p{3.5cm}p{3.5cm}p{3.6cm}p{4.1cm}@{}}
\toprule
 & Provider authorisation with manual capture & On-chain escrow with an oracle & This method \\
\midrule
Who holds the funds & provider & contract account & provider (or contract) \\
Who decides the release & merchant & contract code on oracle data & public rule on signed evidence \\
Recheck by a third party & no, private account history & yes, by re-executing the chain & yes, offline from one file \\
Trust in event reports & merchant's word & oracle operator & named source keys, listed in the report \\
Role of ML & provider-side, opaque & none, or inside the oracle & advisory, logged, outside the rule \\
\bottomrule
\end{tabular}
\end{table}

\section{Evaluation}\label{sec:eval}

\subsection{Risk scorer}\label{sec:risk}

\paragraph{Data and protocol.} PaySim~\cite{lopezrojas2016} simulates 6,362,620 mobile-money transactions over 744
hourly steps, 8,213 of them fraudulent. We split by time: steps 1--400 for training (5,787,030 rows, 4,477 fraud),
401--550 for validation (378,989 rows, 1,588 fraud) and 551--744 for test (196,601 rows, 2,148 fraud). Prevalence
rises from 0.08\% to 0.42\% to 1.09\% across the splits, so the test set is a shifted population and not a random
holdout. Six scale-free features known before confirmation are used (amount relative to balance, account drained,
transaction type, sender history count, amount relative to sender mean, recipient history count). The hour of day
and the absolute amount are dropped: the first is an artefact of the simulator and the second does not transfer
across currencies. Candidates are a class-weighted logistic regression, LightGBM~\cite{ke2017} and a one-line
business rule (an account emptied by a transfer to a recipient never paid before). In
Eqs.~\eqref{eq:select}--\eqref{eq:thr}, $S=5$, $\tau=0.02$, $\rho=0.95$ and $\pi=0.90$. The test set is read once.
Intervals are 95\% percentile intervals from 1,000 stratified bootstrap resamples~\cite{efron1979}.

\paragraph{Deviation.} The first threshold rules (recall $\ge 0.90$ for medium, false-positive rate $\le 0.001$ for
high) produced a high threshold below the medium one. We replaced them with Eq.~\eqref{eq:thr} after the test
results of that first run had been seen. The model ranking does not depend on thresholds, but the precision and
recall figures in Table~\ref{tab:ml} were obtained under rules chosen after one look at the test set, and should be
read with that in mind.

\paragraph{Results.} LightGBM fails the stability constraint: its validation PR-AUC is 0.093, 0.004, 0.815, 0.015
and 0.004 over the five seeds (Figure~\ref{fig:seeds}). With about 4,500 positive rows among 5.8 million and a
minimum leaf size of 200, the trees depend heavily on which rows the bagging draws. An earlier version of the
project had reported a PR-AUC of 0.939 for the same model family, from a two-way split whose thresholds were tuned
on test data; that figure is withdrawn. The logistic regression is selected. On test it reaches a PR-AUC of 0.973
against 0.938 for the business rule; the paired bootstrap difference is 0.035 (95\% CI 0.029--0.041).
Isotonic calibration lowers the Brier score from $3.06\times10^{-3}$ to $1.01\times10^{-3}$ and the expected
calibration error from $1.48\times10^{-2}$ to $8.9\times10^{-4}$. The medium level was set for 95\% recall on
validation and obtains 94.4\% on test (CI 93.4--95.4), so the target is met only at the edge of the interval.

The gap between the rule and the model is small. A single hand-written predicate recovers most of the signal, which
says as much about PaySim as about the models: fraud in the simulator follows one script. We keep the logistic model
because it is calibrated and gives per-feature reasons, not because the benchmark proves it would do well on real
payment traffic.

\begin{table}[t]
\centering\small
\caption{Test results on PaySim steps 551--744 at true prevalence (1.09\%). Thresholds come from validation.
Brackets: 95\% bootstrap intervals. The LightGBM row is the first seed and is shown for reference only.}
\label{tab:ml}
\begin{tabular}{@{}lccccccc@{}}
\toprule
 & & & \multicolumn{2}{c}{Medium level} & \multicolumn{2}{c}{High level} & Val.\ PR-AUC \\
\cmidrule(lr){4-5}\cmidrule(lr){6-7}
Model & ROC-AUC & PR-AUC & Prec. & Recall & Prec. & Recall & sd over seeds \\
\midrule
Logistic regression & 0.9990 & 0.973 {\scriptsize[0.967, 0.978]} & 0.940 & 0.944 & 0.951 & 0.812 & 0.000 \\
Business rule       & 0.9972 & 0.938 {\scriptsize[0.929, 0.948]} & 0.938 & 0.969 & 0.979 & 0.606 & n/a \\
LightGBM (rejected) & 0.9646 & 0.203 & 0.288 & 0.940 & 0.173 & 0.476 & 0.353 \\
\bottomrule
\end{tabular}
\end{table}

\begin{figure}[t]
\centering
\begin{tikzpicture}
\begin{axis}[width=0.62\textwidth, height=4.3cm, font=\footnotesize,
  xlabel={Training seed}, ylabel={Validation PR-AUC}, ymin=0, ymax=1.05, xmin=0.6, xmax=5.4,
  xtick={1,2,3,4,5}, ytick={0,0.25,0.5,0.75,1}, grid=major, grid style={draw=black!10},
  legend style={at={(1.03,0.5)}, anchor=west, draw=none, font=\footnotesize, cells={anchor=west}},
  axis line style={draw=mute}, tick style={draw=mute}]
\addplot[adv, thick, mark=*, mark size=1.8pt] coordinates {(1,0.9523)(2,0.9523)(3,0.9523)(4,0.9523)(5,0.9523)};
\addlegendentry{Logistic regression}
\addplot[cus, thick, dashed, mark=none] coordinates {(0.6,0.9037)(5.4,0.9037)};
\addlegendentry{Business rule}
\addplot[dec, thick, mark=square*, mark size=1.8pt] coordinates {(1,0.0925)(2,0.0042)(3,0.8149)(4,0.0154)(5,0.0042)};
\addlegendentry{LightGBM}
\end{axis}
\end{tikzpicture}
\caption{Validation PR-AUC per training seed. The selection rule of Eq.~\eqref{eq:select} admits a model only if
its standard deviation across seeds is at most 0.02. A single-seed report of LightGBM could have shown 0.81 or
0.004.}
\label{fig:seeds}
\end{figure}

\subsection{Proof packages}

\paragraph{Functional tests.} The verifier test suite has 55 cases, all passing at the evaluated commit. They
cover the paths of Eq.~\eqref{eq:delta}: release, void at the deadline, arbiter hold and ruling, buyer dispute,
evidence signed by an unregistered key, by a key with the wrong role, by a key revoked before or after signing,
evidence for another order or another amount, and back-dated evidence received after the deadline. Two cases run
the JavaScript verifier on the same packages and compare the reports.

\paragraph{Tamper study.} We built three packages (a release under each rule and a void) and, for every file in
each, flipped every bit in turn and ran the verifier with pinned keys: 192,880 trials in total
(Table~\ref{tab:tamper}). Every flip in the terms, the attestations, the receipts, the decision, the key registry
and the manifest was rejected. Two defects appeared, both in the handling of detached signature files.

First, 123 flips were accepted (41 per package), all in \code{manifest.sig.json}. A signature file carries a
\code{created\_at} time next to the signature, and the signature covers the content hash only. The other signature
files are protected because the manifest lists their hashes; the manifest's own signature file is not. Of the 41
accepted flips, 38 turn \code{created\_at} into another well-formed time inside the key's validity interval, and 3
change the unused trailing bits of the base64 signature. None changes what was decided or on what evidence. The
first kind matters all the same: \code{created\_at} feeds the verdict on a key revoked after signing, so an attacker
holding a package signed with a since-revoked operator key could move that time. The fix is to sign the pair
(hash, time) and to decode base64 strictly.

Second, in 1,348 trials the verifier returned no report: the flipped bit made a \code{created\_at} value
unparseable and the date parser raised an exception. A crash is still a refusal, but the verifier's contract is to
report problems and never raise, and a service that calls it without guarding the call could be stopped by a
crafted package. Neither defect is caught by the hand-written test suite. The script
(\code{research/paper/tamper\_experiment.py}) and its raw output are in the repository. Replaying a valid package
takes about 1~ms on one core.

\begin{table}[t]
\centering\small
\caption{Exhaustive single-bit corruption. ``Rejected'': the report says \textsc{invalid}. ``Exception'': the
verifier raised instead of reporting. ``Accepted'': the report says \textsc{valid}; all such flips are in
\code{manifest.sig.json}. Median verification time over 200 runs, one core.}
\label{tab:tamper}
\begin{tabular}{@{}lrrrrrrr@{}}
\toprule
Package & Files & Bytes & Trials & Rejected & Exception & Accepted & Verify (ms) \\
\midrule
\code{delivery-attested/1}, release & 14 & 7,978 & 63,824 & 63,332 & 451 & 41 & 1.05 \\
\code{conditions/1}, release & 16 & 9,788 & 78,304 & 77,736 & 527 & 41 & 1.33 \\
\code{delivery-attested/1}, void & 12 & 6,344 & 50,752 & 50,341 & 370 & 41 & 0.76 \\
\midrule
Total & & & 192,880 & 191,409 & 1,348 & 123 & \\
\bottomrule
\end{tabular}
\end{table}

\subsection{What was not measured}

We did not measure how often the language model drafts the condition a buyer intended; that requires a labelled
set of purchase descriptions and human raters, which we do not have. The carrier, inspector and customs services
are simulations, so no claim is made about integration with real logistics data. Captures were exercised end to
end against a local simulator of the provider's API and the ledger; a full capture on the provider's sandbox is
still outstanding. The scorer was evaluated on synthetic data only.

\section{Limits}\label{sec:limits}

\paragraph{Signatures do not observe the world.} The method reduces ``did the goods arrive'' to ``did the agreed
carrier sign that they arrived''. A dishonest or mistaken source produces a valid package for a false event. The
\code{at\_least} connective lets the parties require agreement among several sources, and the dispute window gives
the buyer a period to object, but neither is a proof of physical fact. The verifier prints this limitation in every
report.

\paragraph{The operator is still a single evaluator.} $\delta$ is public and replayable, so a wrong decision is
detectable after the fact, but it is taken by one service. The operator also issues the receipts that fix $r(a)$,
and could withhold or delay evidence. Moving evaluation of $\delta$ and receipt issuance to the validators would
remove this; the ledger currently anchors outcomes only.

\paragraph{Keys.} A compromised source key can release funds until it is revoked. Validity is checked at reception
time, which limits back-dating but not real-time misuse. Hardware-held keys and short validity intervals are the
standard mitigations, and neither is implemented here.

\paragraph{Models.} Bounding the drafted condition by a schema prevents malformed output from executing. It does
not prevent a well-formed condition that differs from what the buyer meant; the buyer's reading is the only
control. The risk scorer was fitted on a simulator with one fraud pattern and will need re-estimation and drift
monitoring~\cite{dalpozzolo2018} before any use on live traffic.

\section{Conclusion}

Conditional payments can be made checkable by a party who was not involved, without moving custody away from the
provider that already holds the money and without putting a model in the trusted path. The required pieces are
modest: signed evidence with reception receipts, a small decision function that an independent program can replay,
and a rule that model output becomes binding only through terms a human signed. Our tamper study found no undetected change to
terms, evidence or decisions, and two defects in the signature envelope that tests written by hand had missed. Our scoring study found that an unstable model can look
excellent or useless depending on the seed, and that a time-ordered protocol with a stability constraint catches
this before deployment. Two pieces of work come next: replicating the decision function across validators, and a
human-rated study of drafted conditions against buyer intent.

\paragraph{Reproducibility.} All paths are relative to the repository root. The model evaluation script is
\code{guardian/\allowbreak training/\allowbreak evaluate.py} and its raw output is under
\code{research/\allowbreak results/\allowbreak raw/\allowbreak ml-eval/}. The PaySim file is identified by SHA-256
\code{16910f90\ldots4eea6b}. The verifier tests are in \code{verifier/tests} and the tamper study in
\code{research/paper/}.

{\small
\begin{thebibliography}{99}
\setlength{\itemsep}{1pt}

\bibitem{asokan2000} N.~Asokan, V.~Shoup, and M.~Waidner. Optimistic fair exchange of digital signatures.
\emph{IEEE Journal on Selected Areas in Communications}, 18(4):593--610, 2000.

\bibitem{bartoletti2018} M.~Bartoletti and R.~Zunino. BitML: A calculus for Bitcoin smart contracts. In
\emph{Proc.\ ACM SIGSAC Conference on Computer and Communications Security (CCS)}, pages 83--100, 2018.
doi:10.1145/3243734.3243795.

\bibitem{bernstein2012} D.~J. Bernstein, N.~Duif, T.~Lange, P.~Schwabe, and B.-Y. Yang. High-speed high-security
signatures. \emph{Journal of Cryptographic Engineering}, 2(2):77--89, 2012. doi:10.1007/s13389-012-0027-1.

\bibitem{buchman2018} E.~Buchman, J.~Kwon, and Z.~Milosevic. The latest gossip on BFT consensus. arXiv:1807.04938,
2018.

\bibitem{caldarelli2020} G.~Caldarelli. Understanding the blockchain oracle problem: A call for action.
\emph{Information}, 11(11):509, 2020. doi:10.3390/info11110509.

\bibitem{dalpozzolo2018} A.~Dal~Pozzolo, G.~Boracchi, O.~Caelen, C.~Alippi, and G.~Bontempi. Credit card fraud
detection: A realistic modeling and a novel learning strategy. \emph{IEEE Transactions on Neural Networks and
Learning Systems}, 29(8), 2018. doi:10.1109/TNNLS.2017.2736643.

\bibitem{davis2006} J.~Davis and M.~Goadrich. The relationship between precision-recall and ROC curves. In
\emph{Proc.\ 23rd International Conference on Machine Learning (ICML)}, pages 233--240, 2006.
doi:10.1145/1143844.1143874.

\bibitem{efron1979} B.~Efron. Bootstrap methods: Another look at the jackknife. \emph{The Annals of Statistics},
7(1):1--26, 1979. doi:10.1214/aos/1176344552.

\bibitem{greshake2023} K.~Greshake, S.~Abdelnabi, S.~Mishra, C.~Endres, T.~Holz, and M.~Fritz. Not what you've
signed up for: Compromising real-world LLM-integrated applications with indirect prompt injection. In \emph{Proc.\
16th ACM Workshop on Artificial Intelligence and Security (AISec)}, pages 79--90, 2023. doi:10.1145/3605764.3623985.

\bibitem{haber1991} S.~Haber and W.~S. Stornetta. How to time-stamp a digital document. \emph{Journal of
Cryptology}, 3(2):99--111, 1991. doi:10.1007/BF00196791.

\bibitem{rfc8032} S.~Josefsson and I.~Liusvaara. Edwards-curve digital signature algorithm (EdDSA). RFC~8032, IETF,
2017.

\bibitem{kapoor2023} S.~Kapoor and A.~Narayanan. Leakage and the reproducibility crisis in machine-learning-based
science. \emph{Patterns}, 4(9):100804, 2023. doi:10.1016/j.patter.2023.100804.

\bibitem{ke2017} G.~Ke, Q.~Meng, T.~Finley, T.~Wang, W.~Chen, W.~Ma, Q.~Ye, and T.-Y. Liu. LightGBM: A highly
efficient gradient boosting decision tree. In \emph{Advances in Neural Information Processing Systems 30}, 2017.

\bibitem{lamela2018} P.~Lamela~Seijas and S.~Thompson. Marlowe: Financial contracts on blockchain. In
\emph{Leveraging Applications of Formal Methods, Verification and Validation (ISoLA)}, LNCS 11247. Springer, 2018.
doi:10.1007/978-3-030-03427-6\_27.

\bibitem{rfc6962} B.~Laurie, A.~Langley, and E.~Kasper. Certificate transparency. RFC~6962, IETF, 2013.

\bibitem{lopezrojas2016} E.~A. Lopez-Rojas, A.~Elmir, and S.~Axelsson. PaySim: A financial mobile money simulator
for fraud detection. In \emph{Proc.\ 28th European Modeling and Simulation Symposium (EMSS)}, Larnaca, 2016.

\bibitem{lundberg2017} S.~M. Lundberg and S.-I. Lee. A unified approach to interpreting model predictions. In
\emph{Advances in Neural Information Processing Systems 30}, pages 4765--4774, 2017.

\bibitem{niculescu2005} A.~Niculescu-Mizil and R.~Caruana. Predicting good probabilities with supervised learning.
In \emph{Proc.\ 22nd International Conference on Machine Learning (ICML)}, pages 625--632, 2005.
doi:10.1145/1102351.1102430.

\bibitem{nosek2018} B.~A. Nosek, C.~R. Ebersole, A.~C. DeHaven, and D.~T. Mellor. The preregistration revolution.
\emph{Proceedings of the National Academy of Sciences}, 115(11):2600--2606, 2018. doi:10.1073/pnas.1708274114.

\bibitem{pagnia1999} H.~Pagnia and F.~C. G\"artner. On the impossibility of fair exchange without a trusted third
party. Technical Report TUD-BS-1999-02, Darmstadt University of Technology, 1999.

\bibitem{peytonjones2000} S.~Peyton~Jones, J.-M. Eber, and J.~Seward. Composing contracts: An adventure in
financial engineering. In \emph{Proc.\ 5th ACM SIGPLAN International Conference on Functional Programming (ICFP)},
pages 280--292, 2000. doi:10.1145/351240.351267.

\bibitem{rfc8785} A.~Rundgren, B.~Jordan, and S.~Erdtman. JSON canonicalization scheme (JCS). RFC~8785, IETF,
2020.

\bibitem{saito2015} T.~Saito and M.~Rehmsmeier. The precision-recall plot is more informative than the ROC plot
when evaluating binary classifiers on imbalanced datasets. \emph{PLoS ONE}, 10(3):e0118432, 2015.
doi:10.1371/journal.pone.0118432.

\bibitem{szabo1997} N.~Szabo. Formalizing and securing relationships on public networks. \emph{First Monday},
2(9), 1997. doi:10.5210/fm.v2i9.548.

\bibitem{zadrozny2002} B.~Zadrozny and C.~Elkan. Transforming classifier scores into accurate multiclass
probability estimates. In \emph{Proc.\ 8th ACM SIGKDD International Conference on Knowledge Discovery and Data
Mining (KDD)}, 2002. doi:10.1145/775047.775151.

\bibitem{zhang2016} F.~Zhang, E.~Cecchetti, K.~Croman, A.~Juels, and E.~Shi. Town Crier: An authenticated data
feed for smart contracts. In \emph{Proc.\ ACM SIGSAC Conference on Computer and Communications Security (CCS)},
pages 270--282, 2016. doi:10.1145/2976749.2978326.

\bibitem{zhang2020} F.~Zhang, D.~Maram, H.~Malvai, S.~Goldfeder, and A.~Juels. DECO: Liberating web data using
decentralized oracles for TLS. In \emph{Proc.\ ACM SIGSAC Conference on Computer and Communications Security
(CCS)}, pages 1919--1938, 2020. doi:10.1145/3372297.3417239.

\end{thebibliography}
}
\end{document}
